\section{Conclusion}
\label{sec:conclusion}

In this paper, we introduce a tractable second-order cone formulation of the harmonic hosting capacity problem for balanced three-phase networks which uses harmonic current injections in rectangular coordinates for all harmonics in accordance with the existing power quality standards. The problem formulation allows to determine a lower bound on the harmonic hosting capacity under the assumption of equal phase angles of all harmonic injections, which does not allow for harmonic compensation. It is numerically demonstrated that the second-order cone formulation is equivalent to the nonlinear formulation of the problem using an industrial network as a test case.  Furthermore, the numerical results obtained the 1888 bus Rte test system show that this approach can successfully be applied to realistic country-scale networks.  

It is further shown that the choice of the fairness criterion for allocating the harmonic budget results in substantially different system harmonic hosting capacities which should be addressed in future editions of power quality standards. The allocation of the harmonic budget will gain evermore importance with the massive expected integration of inverter-based resources in the future.

Furthermore, although particular fairness criteria result in substantially different harmonic spectra on the nodal level, the average harmonic voltages on the system level behave similarly despite having different individual magnitudes.

The presented models do not yet capture the stochastic nature of injections or uncertainties around grid parameters affecting the harmonic hosting capacity of a power system. In future work, the authors will work towards including both the stochastic nature of the power system itself, e.g., component parameters or power system exploitation, as well as the angle of the harmonic current injected by the units.